\cotefile{83}
\scanpage{162}
% L1–6
La \deleted{catégorie} donnée d'une var. quadratique
$X$ de $S$, revient à la donnée d'un
\deleted{\unsure{tel syst.|fibré inv.}} $L$ affine et d'un div. quadr. \unsure{dans|\deleted{dans}} $\deleted{\unsure{V}}\check{\mathbb V}(L)$
(lequel sera $X$), ou encore à $L$ et
la fonction de degré $\leq 2$ sur $\check{\mathbb V}(L)$,
dont la partie homogène soit $u\mapsto u^{\otimes2}$.

% L7–13
\deleted{(Dualement, cela \unsure{revient|revi}}
L'algèbre de $X$ \deleted{\unsure{sera|dans}} le dual \inserted{$H\simeq E\otimes L^{-1}$} de $E$,
l'élém$^{\mathrm t}$ unité étant le dual de l'augmentation $E\to k$. \unsure{L$^{k}$a|La} \unsure{transformée|transformation}
par $\otimes L^{-1}$ de $q$ sera une forme
quadratique sur $H$, (\unsure{en|ce} sera la forme
déterminant) égale à $1$ sur $\mathbf{1}$. Il
% L14–21
y aura sur $E$ une unique structure
\deleted{multiplicative} d'alg. quadratique, pour
laquelle $1$ est unité et $q_H$ est le
déterminant. On récupère $\mathrm{Tr}$
\deleted{comme} via la forme bil. associée à $q_H$
comme $x\mapsto\varphi_H(1,x)$, et dual$^{\mathrm{mt}}$ sur
$L$, on récupère $\mathbf{t}\in E$, \deleted{comme} par
la formule
\[
\varphi_{q_E}\bigl(\underset{\substack{\rotatebox{-90}{$\in$}\\L}}{\unsure{u}},\underset{\substack{\rotatebox{-90}{$\in$}\\E}}{x}\bigr)
=\deleted{\unsure{x}}\underset{\substack{\rotatebox{-90}{$\in$}\\L}}{u}\otimes
\underbrace{(\deleted{\unsure{u}}x\wedge\mathbf{t})}_{\in\,\det E\simeq L}\in\deleted{\det E}=L^{\otimes 2}
\]
% L21 lower left
\margin{Attention aux conv. de signe!}
% archive sheet number: 161
\scanpage{163}
% L1–5
Ceci \inserted{(i.e. $\mathbf t$)} définit une involution sur
$E$ compatible avec l'augmentation, i.e.
une involution sur $A$ i.e. sur $\check{\mathbb V}(A)$,
laquelle laisse le diviseur $X$ invariant --
et y induit l'involution canonique.

% L6–19
Attention, si on se donne d'\uline{avance}
une involution sur $L$, et cherche les
div. quadr. sur $\check{\mathbb V}(L)$ qui \deleted{\unsure{soient stables|sont stables}} \inserted{la redonnent,}
\inserted{celles-ci sont stables} par l'involution, \deleted{ceci soit par \unsure{né.}}
mais cette condition n'est pas suffisante
-- en car. $2$, on peut prendre l'involution
identique, alors qu'il y aura des
diviseurs quadratiques (scindés, p.ex)
pour lesquels l'involution n'\uline{est} \uline{pas}
l'identité! Mais \deleted{\unsure{si|car}} $2$ inv., cette difficulté
disparaît -- cela se ramène à voir
que tt div. sur $\check{\mathbb V}(\deleted{\unsure{A}}L)\deleted{\unsure{E'|E}}$, stable
par $\lambda\to-\lambda$, est bien l'image inverse
d'un $\deleted{\unsure{\zeta}}\unsure{\zeta|\xi}\in L^2$ par $L\xrightarrow{u\mapsto u^{\otimes 2}} L^{\otimes 2}$.

% L20–24
(Il suffit $=$ que $2$ ne soit pas
div. de zéro \unsure{\ldots} \deleted{\unsure{?}} Il paraît que cette
dernière propriété soit vraie -- mais
pour l'appliquer aux $A$ affines \unsure{quelc.|quelq.}
il faut supposer que l'involution sur $A$
a un p$^{\mathrm t}$ fixe\ldots)
% archive sheet number: 162
\scanpage{164}
% L1–3
donc l'application
\[
\zeta\longmapsto q_\zeta,\quad A*A\longrightarrow
\begin{aligned}&\text{formes quadratiques}\\&\text{sur $E$ qui sont}\\&\text{compatibles à $t$}\end{aligned}
\]
est compatible avec l'application
\[
\begin{aligned}
\deleted{A}L^{\otimes 2}&\longrightarrow L^{\otimes 2}\\
\alpha&\longmapsto-\alpha
\end{aligned}
\]
% L4–7
sur les \deleted{\unsure{modules d'opérateurs|modules d'opération}} \inserted{modules des translations}. Donc il
vaut mieux changer notre convention
et \deleted{\unsure{poser|pour}} prendre $q_\zeta$ tel que
\[
\zeta+q_\zeta(x)=x*x
\]

% L8–10
Ainsi, on a \unsure{un iso.|un isom.} de torseurs
\[
\mathrm{Quadcomp}\underbrace{(E,\epsilon,t)}_{\text{ou }(L,\sigma)}\simeq L*_{\sigma}L
\]

% L11–16
Moyennant quoi \unsure{on} va \unsure{identifier|identifie} les
deux. Ceci posé, pour toute $q\in L*_\sigma L$,
le diviseur quadratique associé n'est
autre que l'image inverse de la section
$q$ de $\check{\mathbb V}(L*_\sigma L)$ par $\mathbb V(L)\mathbin{\text{---}}\mathbb V(L*_\sigma L)$.

% L17–21
\S En termes de coordonnées,
i.e. choisissant une origine dans $\deleted{\unsure{L}}A$,
de sorte que $A$ s'identifie à $L$,
et $t$ est \deleted{\unsure{une}} donné par $\tau\in L$,
\[
\sigma_L(\deleted{\unsure{k}}x)=\tau-x
\]
% archive sheet number: 163
\scanpage{165}
% L1–9
Cherchons les formes quadratiques \unsure{(E}
sur $E$ compatibles avec une involution déjà donnée, \inserted{i.e. avec un $t\in E$ donné.}
Il est immédiat \unsure{(.|(}
que ces formes forment un torseur
sous l'\unsure{ens|un} des formes \inserted{(à val. dans $L^{\otimes 2}$)}
qui se factorisent par $E/L\simeq k$, i.e. sous
$\deleted{\unsure{k}}L^{\otimes 2}$ lui $=$.
On en a déjà un tel torseur, savoir $A*A$, et une
application quadratique \unsure{canonique|homogène}
\[
A\deleted{*A}\longrightarrow A*A
\]
% L11–14
(``compatible avec chg$^{\mathrm t}$ de base'' donc
\unsure{défini|defini|définie} sur les schémas correspondants)
savoir
\[
x\,\unsure{\mathrel{\vdash}|\mapsto}\,x*x
\]

% L15–22
Ainsi, \deleted{\unsure{soit}} pour \deleted{tout}
\[
\zeta\in A*A,
\]
si on dénote par $q_\zeta$ l'unique \deleted{forme} \inserted{fonction}
\deleted{quadratique} \inserted{polyn. de degré $\leq 2$}
\[
q_\zeta:\deleted{\unsure{E}}A\longrightarrow L^{\otimes 2}
\]
qui satisfait
\[
\zeta\deleted{\unsure{-}}+q_\zeta(x)=x*x\qquad \forall\,x\text{ dans }A
\]
on aura, pour $\alpha\in L^{\otimes 2}$, \unsure{remplaçant}
$\zeta$ par $\zeta+\alpha$
\[
(\zeta+\alpha)+q_{\zeta+\alpha}(\unsure{u|x})=x*x
\]
donc
\[
q_{\zeta+\alpha}(x)\deleted{\unsure{+}}-q_\zeta(x)=
\underset{\substack{\Vert\\-\alpha\epsilon(x)^2}}{\deleted{\unsure{+}}-\alpha}
\quad\begin{aligned}&\text{pour }x\in A\\&\text{i.e. }\epsilon(x)=1\end{aligned}
\]
% archive sheet number: 164
\scanpage{166}
% L1–10
On peut construire un torseur
$Q_{\sigma}\deleted{\unsure{(L)}}=Q(L,\sigma)$ sous $L^{\otimes 2}$, et une fonction
\unsure{para-quadratique|sous-quadratique}
\[
A\longrightarrow Q
\]
dont la partie homogène soit
\[
\unsure{x|u}\,\unsure{\vdash|\mapsto}\,u^{\otimes 2}:L\mathbin{\text{---}}L^{\otimes 2},
\]
et qui soit \uline{universelle} pour ces
propriétés, \deleted{Pour ceci, on note que}
\inserted{de sorte que les \deleted{fonctions}}
\unsure{ft-quadratiques|f.-quadratiques|sous-quadratiques} à val. dans $L^{\otimes 2}$ comp. avec
$\sigma$ soient simplement les composées
\diagram{\begin{tikzcd}[ampersand replacement=\&,column sep=large]
 A\arrow[r]\& Q\arrow[r, "{\substack{\mathrm{iso.}\ \mathrm{de}\\L^{\otimes 2}\text{-tors.}\\(\text{\unsure{correspond$^{\mathrm t}$|correspondant}}\\\text{au choix d'une}\\\text{origine})}}"']\&L^{\otimes 2}
\end{tikzcd}}

% L12–18
Pour ceci on \unsure{note|voit} que l'ens des
fonctions \deleted{\unsure{sur|sous}} quadr. (à valeurs dans $L^{\otimes 2}$)
qui nous conviennent, forment un
torseur sous le module des formes
quadr. qui se factorisent par
$E/\unsure{E^0|E_0|E^\circ}\simeq k$, lequel module est \unsure{isom|iso}
à $\mathrm{Quad}(k,L^{\otimes 2})\simeq L^{\otimes 2}$. En termes
% archive sheet number: 165
\scanpage{167}
% L1–8
de fonctions \unsure{semi-quadratiques|sous-quadratiques} sur $A$,
cela signifie simplement qu'elles
forment un torseur sous $L^{\otimes 2}$, par
addition de fonctions \uline{constantes}.
Si $\Lambda$ est ce torseur, on a
\[
\begin{aligned}
\Lambda\times A&\mathbin{\text{---}}\deleted{\unsure{A}}L^{\otimes 2}\\
(q,x)&\mathrel{\vdash\!\!-\!\!-}q(x)
\end{aligned}
\]
tel que
\[
(q+\alpha,x)\,\unsure{\mathrel{\vdash\!\!-\!\!-}|\mapsto}\,q(x)+\alpha\qquad\alpha\in L^{\otimes 2}
\]

% L9–13
et on prend
\[
Q_\sigma=\mathrm{Hom}_{\substack{\text{torseurs }\deleted{\text{su}}\\\text{sur }L^{\otimes 2}}}(\Lambda,L^{\otimes 2})
\]
% slanted
\margin{torseur inverse de $\Lambda$}
d'où
\[
\Lambda\simeq\mathrm{Hom}_{\text{tors sur }L^{\otimes 2}}(Q,L^{\otimes2})
\]
d'où \deleted{\unsure{accouplement|canoniquement}} application can.
\[
A\longrightarrow Q_\sigma
\]
% L13 left bubble
\margin{C'est $x\mapsto x\deleted{\unsure{\otimes}}\unsure{\cdot|\otimes}\sigma x\in\mathrm{Sym}^2(E)
\simeq\mathrm{Quad}(\underset{\substack{\Vert\\ E\otimes L^{-1}}}{\check E},k)
\simeq\mathrm{Quad}(E,L^{\otimes 2})$}

% L14–21
\deleted{Pour un choix}
Le choix d'une $q$ \inserted{$\in\Lambda$} revient à celui
d'un \unsure{isom|section} de $L^{\otimes 2}$-torseurs i.e. d'une
origine dans $Q_\sigma$ (car $Q_\sigma\simeq\Lambda$,
mais cette bijection renverse l'op.
de $L^{\otimes 2}$ en la symétrique $-$) Ceci dit,
le diviseur quadratique \inserted{$X_{\unsure{\zeta|q}}$} associé à une
$q$, correspondant à une section $q$ de $Q_\sigma$,
% archive sheet number: 166
\scanpage{168}
% L1–2
n'est autre que l'image inverse de cette
section.

% L3–12
\deleted{\unsure{L|On}}On voit que la bonne notion ici
est celle de fonctions \unsure{semi|sous} quadratiques
d'un torseur $A$ dans un autre
$Q$, sous $L^{\otimes 2}$ cette fois, ``compatible''
avec \deleted{\unsure{$\sigma$}} ($u\mathrel{\vdash\!\!-\!\!-}u^{\unsure{2}}$ de $L$ dans $L^{\otimes 2}$ et) $t$: ce
sont les fonctions \inserted{\unsure{para|sous}} quadratiques
\deleted{\unsure{$q$}} \deleted{telles}.
\[
q:A\longrightarrow Q
\]
telles que
\[
q(x+\underset{\substack{\rotatebox{-90}{$\in$}\\L}}{\alpha})=q(x)-
\underset{\substack{\rotatebox{-90}{$\in$}\\L}}{(t\wedge x)}\otimes\alpha+
\alpha\deleted{\unsure{\otimes}}\otimes\alpha
\]
% slanted
\margin{attention conv. de signe!}

% L13–17
Deux telles fonctions \inserted{\unsure{para|sous}} quadratiques diffèrent
par une constante. Si on en a une,
cela fixe $Q$ à iso. près comme
sol. du pb universel. Et on vient
d'en construire une tantôt!

% L18–20
\deleted{On peut}
Ainsi, à \unsure{tte|tt} ``droite affine'' $A$ à involution
on a associé $Q(A,\sigma)=Q(A)$
et $A\longrightarrow Q(A)$ rev. de degré $2$
% archive sheet number: 167
\scanpage{169}
% L1–7
Considérons maintenant
\[
\underset{\text{sous }L_1}{(A_1,\sigma_1)}\quad
\underset{\text{sous }L_2}{(A_2,\sigma_{\unsure{2}})}
\]
droites affines à involution, d'où
\[
\underset{\text{\unsure{sous|sur} }L_1\otimes L_2}{A=\deleted{Q(A_1}A_1*A_2}
\quad\deleted{\unsure{b_1}}
\]
je dis que
\[
\underbrace{Q(A)}_{\text{sous }L^{\otimes 2}}\simeq
\underbrace{Q(A_1)}_{\text{sous }L_1^{\otimes 2}}
\deleted{\otimes}{*}\underbrace{\deleted{\unsure{A}}Q(A_2)}_{\text{sous }L_2^{\otimes 2}}
\]
\margin{??}
(iso canonique)

% L8–10
\margin{?}Pour donner un sens à cette
formule, il faut définir sur un
$Q(A)$ une involution naturelle

% L11–17
En tous cas, \deleted{en termes} on a
une application canonique
\[
\begin{aligned}
&q_1,q_2\mathbin{\text{---}}q_1\otimes q_2\mid \unsure{\check A_1|\dot A_1|A_1}*A_2\subset E_1\otimes E_2\\
&Q(A_1)\times Q(A_2)\longrightarrow Q(A)
\end{aligned}
\]
\deleted{\unsure{donnée en termes de co}}
qui s'explicite ainsi, si on choisit
\deleted{\unsure{pour|base}} pt base $o_1\in A_1$, $o_2\in A_2$, de sorte
que
\[
A_i\simeq A_i^0,\quad \sigma_i\text{ devient }x\mathrel{\vdash\!\!-\!\!-}b_i-x
\]
% L18–21
les formes quadratiques sur $A_i$ \deleted{\unsure{dev|comp}}compatibles
à $\sigma_i$ sont
\[
q_{c_i}\text{ ou }q_{b_i,c_i}\quad x_i\longmapsto x_i^2-\deleted{\unsure{t_i}}b_i\unsure{x.|x_i}+c_i
\qquad c_i\in L^{\otimes 2}\quad\unsure{avec|conv.|comme}
\]
et on a alors
\[
q_{c_{\unsure{i|1}}}\otimes q_{\unsure{c_2|b_2}}=q_c\quad
\underset{\substack{\rotatebox{-90}{$\in$}\\L_1\otimes L_2}}{x}
\mathrel{\vdash\!\!-\!\!-}x_{\deleted{\unsure{12}}}^{2}-(b_1b_2)x_{\unsure{2|{}}}+
(b_2^2c_1+b_1^2c_2-4c_1c_2)
\]
% archive sheet number: 168
\scanpage{170}
% L1–4
i.e
\[
q_{c_1}\otimes q_{c_2}=q_c\qquad c=b_2^2c_1+b_1^2c_2-4c_1c_2
\]
Cela montre que pour $\alpha_1\in L_1^{\otimes 2}$, on a
\[
(q_{\unsure{4|1}}\deleted{\unsure{\oplus}}+\alpha)\otimes q_2=q_1\otimes q_2+
\underset{\substack{\rotatebox{-90}{$\in$}\\L_1^{\otimes 2}}}{\alpha}
\underset{\substack{\rotatebox{-90}{$\in$}\\L_2^{\otimes 2}}}{\delta_2(q_2)}
\]
% L5–9
(il suffit de faire $c\mapsto c+\alpha$ dans la formule
pour $c$) et de $=$ 
\[
q_1\otimes(q_2+\beta)=q_1\otimes q_2+\delta_1(q_1)\beta
\]
où
\[
\begin{aligned}
\delta_1(q_{c_1})&=b_1^2-4c_{\unsure{i|1}}\\
\delta_2(q_{c_2})&=b_2^2-4c_2 .
\end{aligned}
\]

% L10–12
\deleted{Ainsi, il faut regarder}
De plus
\[
\underset{\substack{\vert\\\text{discriminant}\\\text{de formes sur}\\\text{module \unsure{env.}}\\\text{de }A_1*A_2=A\\\rotatebox{-90}{$\in$}\\(L_1\otimes L_2)^{\otimes 2}}}{\delta(q_1\otimes q_2)}
=\underset{\substack{\rotatebox{-90}{$\in$}\\L_1^{\otimes 2}}}{\delta(q_1)}
\underset{\substack{\rotatebox{-90}{$\in$}\\L_2^{\otimes 2}}}{\delta(q_2)}
\]

% L13–17
On trouve diagramme commutatif
\diagram{\begin{tikzcd}[ampersand replacement=\&,row sep=large,column sep=huge]
 A_1\times A_2\arrow[r,"*"]\arrow[d,"\pi_1\times\pi_2"] \& A=A_1*A_2\arrow[d,"\pi_A"]\\
 Q_1\times Q_2\arrow[r,"\otimes"]\arrow[d,"\delta_{A_1}\times\delta_{A_2}"] \& Q=Q_{A,\sigma_A}\arrow[d,"\delta_A"]\\
 L_1^{\otimes 2}\times L_2^{\otimes 2}\arrow[r,"{\overset{\mathrm{can}}{\sim}}"] \& (L_1\unsure{\otimes|\cdot|{}}L_2)^{\otimes 2}
\end{tikzcd}}
où, en termes des \deleted{\unsure{coor}} choix des pt$^{\mathrm s}$ marqués
$o_1$ dans $A_1$, $o_2$ dans $A_2$, $\unsure{\sigma_1|o_1}*o_2$ dans $A=A_1\unsure{\times|*|\alpha}A_2$,
tout est explicité (identifiant $A_1,A_2,A_1*A_2$ à
$L_1,L_2$ et $L_1\otimes L_2$ resp.):
% archive sheet number: 169
\scanpage{171}
% L1–4
\[
\begin{aligned}
x_1*x_2&=(o_1+x_1)*(\unsure{o|o_2}+x_2)
=o_1*o_2+\bigl(x_1\otimes(\overbrace{\sigma_2(o_2)}^{b_2}-o_2\bigr)\\
&\qquad+(\sigma_1(o_1)-o_1)\otimes x_2-2\deleted{\unsure{f}}x_1\otimes x_2\\
&=o+\underbrace{x_1\otimes b_2+b_1\otimes x_2-2x_1x_2}\\
&=x_1b_2+b_1x_2-2x_1x_2
\end{aligned}
\]
% L5–10
$\pi_1:A_1\longrightarrow Q$ transforme $x_1$ en la forme
\[
\pi_1(x_1)=q_{x_1(b-x_1)}
=\bigl(z_1\mapsto {\deleted{\unsure{x}}z}_1^2-b_1z_1+
\deleted{\unsure{(}}x_1(b_1-x_1)\unsure{'}=(z_1-x_1)(z_1-(b_1-x_1))\bigr)
\]
et de $=$ pour $\pi_2$, donc
$\pi_1(x_1)\otimes\pi_2(x_2)$ est la forme
\[
\begin{aligned}
z\longmapsto{}&z^2-b_1b_2z+b_2^2x_1(b_1-x_1)+b_1^2x_2(b_2-x_2)\\
&-4x_1(b_1-x_1)x_2(b_2-x_2)
\end{aligned}
\]
% L11–12
et
\[
\pi(x_1*x_2)=\bigl(z\mapsto z^2-\underbrace{b_1b_2}_{b}z+
\underbrace{(x_1*x_2)(b_1b_2-x_1*x_2)}_{\substack{(x_1b_2+b_1x_2-2x_1x_2)(b_1b_2-x_1b_2\unsure{+|-}b_1x_2\\{}+2x_1x_2)}}
\]

% L13–19
Le calcul de l'égalité est un peu pénible
mais ça peut se faire (j'ai dû le
vérifier précéd$^{\mathrm{mt}}$ par un autre biais,
en choisissant $q_1\,q_2$ et faisant la vérification
en choisissant $o_i=x_i$ \deleted{tel que $q_{\unsure{i|1}}(o_i)=0$} i.e
il suffit de le faire pour $\unsure{x_1'|x_1}=0$, $\unsure{\alpha_2|x_2}=0$
où \unsure{c'est|est} trivial.

% L20–26
On voudrait écrire \deleted{$Q(A)=Q(A$}
\[
Q(A_1*A_2)\simeq Q(A_1)*Q(A_2)
\]
mais la théorie du $*$ développée
précédemment ne s'y prête pas, \unsure{car},
\deleted{$Q$} les $Q(A_i)$ ne sont pas involutifs.
Il faut reprendre \deleted{tout} le \unsure{formalisme|formulation|formalisation} \inserted{\unsure{$*$|$A$|($A$}}
de façon plus générale!
% archive sheet number: 170
\scanpage{172}
% L1–2
Considérons une extension
\diagram{\begin{tikzcd}[ampersand replacement=\&,column sep=small]
 \unsure{\cdot|0}\arrow[r,dash]\&E^0\arrow[r,dash]\&E\arrow[r,"\epsilon"]\&
 \underset{\substack{\rotatebox{-90}{$\simeq$}\\ E/E^0}}{F}\arrow[r,dash]\&0
\end{tikzcd}}
% L3–10
nous allons déterminer les involutions
de $E$, induisant $-\mathrm{id}_{E^0}$ sur $\unsure{E^0|E_0^0}$, et
$+\mathrm{id}_F$ dans $F$. Un endom \inserted{$\sigma$} de $E$ ayant
cette propriété (involutif ou pas) doit
être tel que
\[
u(x)=x+\sigma(x)=0\quad\text{si }x\in E^0,
\]
\margin{(exprime que $\sigma\mid E^0=-\mathrm{id}_{E^0}$}
donc $u$ se factorise en
\[
\pi:F\longrightarrow E\qquad u=\pi\cdot\epsilon
\]
% L11–12
et \deleted{$u(x)$} $\sigma$ se récupère via $\pi$ comme
\[
\deleted{\unsure{u(x|u(x)}} \sigma(x)=\pi(\epsilon(x))-x
\]

% L13–19
Si on exprime maintenant que
$\sigma$ induit $\mathrm{id}_F$ dans $E/E^0$, on trouve
la condition \(\displaystyle x\equiv\pi\epsilon(x))-x\quad(E^0)\quad\forall \unsure{1|x}\in E\)
i.e. \(\displaystyle \pi\underbrace{\epsilon(x)}_{y}\equiv2x\qquad\forall x\in E\)
ce qui signifie aussi que \unsure{le|les|la} $\epsilon$ des deux
\unsure{mbres|membres} est le $=$, soit (posant $\deleted{\unsure{\pi}}\epsilon(x)=y$
\[
\boxed{\epsilon\pi(y)=2y}
\]

% L20–22
Ainsi, $\pi$ est ``presque'' une section
c'est une $2$-section. \deleted{Cette} Plus \unsure{généralt|généralement},
si $\lambda\in k$ et si on a une application
% archive sheet number: 171
\scanpage{173}
% L1–8
$\pi:F\mathbin{\text{---}}E$ telle que
\[
\epsilon\pi=\lambda\,\mathrm{id}_F,
\]
cela signifie que l'extension
$ {}^{\lambda}E$ de $F$ par $E$, déduite de $E$ par
image inverse par l'hom. $\lambda\,\mathrm{id}_F$ \unsure{de|dans} $F$,
ou par image directe par $\lambda\,\mathrm{id}_E$ de $E$,
est munie d'un scindage.
On dit que $E$ est $\lambda$-scindée.

% L9–14
Si p.ex. $F=k$ Cas \uline{d'}\uline{un} espace
affine $A=\epsilon^{-1}(1)$ sous $E^0$), la
donnée de $\pi$ \inserted{$:k\unsure{\mathbin{\text{---}}|\longrightarrow} E$} revient à la donnée
d'un élément $t$ de $E$, et la
condition envisagée c'est que
\[
\epsilon(t)=\lambda .
\]

% L15–21
Cela signifie aussi que le torseur
\[
\epsilon^{-1}(\deleted{\unsure{1|A}}\lambda)={}^{\lambda}A
\]
déduit de $A$ par l'homothétie $\lambda\,\mathrm{id}_{A^0=E^0}$
sur le module des translations, est
muni d'une section \inserted{``origine''}. Donc
l'application
\[
\begin{aligned}
x&\longmapsto\lambda x && A\mathbin{\text{---}}{}^\lambda A\\
&\text{compatible avec} && A^0\mathrel{\mathop{\text{---}}\limits^{\lambda\,\mathrm{id}}} A^0
\end{aligned}
\]
% archive sheet number: 172
\scanpage{174}
% L1–8
peut être interprétée, par l'identification
\[
\begin{aligned}
{}^\lambda A&\xrightarrow{\sim}A^0\\
x&\mapsto x-t
\end{aligned}
\]
comme une application de $A$ dans $A^0$,
de façon précise on définit
\[
u:A\longrightarrow A^0\qquad u(x)=\lambda x-t
\]
application affine compatible avec
\[
A^0\mathrel{\mathop{\text{---}}\limits^{\lambda\,\mathrm{id}_{A^0}}}A^0
\]
% L9–10, between lines
\inserted{et inversement, une telle application s'interprète comme
\deleted{\unsure{application|applic. en}} \unsure{isomorphisme|iso} de $A^0$-torseurs $ {}^\lambda A\simeq A^0$, i.e. choix d'une origine \unsure{$*$}}
\margin{$*$ de $ {}^\lambda A$}

% L10–15
Dans le cas où $\lambda=2$, donc \unsure{où $t$ correspond}
à une involution de $F$ au sens
du début de ces notes, on a donc
\[
u(x)=2x-t=-(\sigma(x)-x),
\]
c'est donc au signe près notre fameux
$\sigma(x)-x$.

% L16–22
Or dans notre construction \unsure{de|des}
$*A_i$, ce sont ces applications $u_i(x_i)$
(ou ce qui revient au $=$, au signe
près, les $\sigma_i(x_i)-x_i$) qui jouaient le
rôle crucial. \deleted{Il n'y avait pas}
Cette application $u(x)$ étant affine,
se prolonge linéairement en
\[
E\xrightarrow{u}A^0\qquad
u(x)\unsure{\mathrel{=\!\!-}|=}\lambda x-\underbrace{\epsilon(x)t}_{=\pi\epsilon(x)}
\]
% archive sheet number: 173
\scanpage{175}
% L1–7
Revenant au cas d'une ext. générale
\diagram{\begin{tikzcd}[ampersand replacement=\&,column sep=small]
0\arrow[r,dash]\&E^0\arrow[r,dash]\&E\arrow[r,dash]\&\underset{\substack{\Vert\\E/E^0}}{F}\arrow[r]\&0
\end{tikzcd}}
et d'une $\lambda$-section
\[
\pi:F\longrightarrow E,
\]
associons lui
\[
\begin{aligned}
\deleted{\pi^{\unsure{t|*|\sharp|\pi}}}u\unsure{:=|:}&\,E\mathbin{\text{---}}E^0\\
\deleted{\pi^{\unsure{t|*|\sharp|\pi}}}u(x)&=\lambda x-\pi\epsilon(x)
\end{aligned}
\]
% L8–11
Elle envoie bien $E$ dans $E^0$, car
\[
\epsilon\deleted{\pi}u(x)=\lambda\underbrace{\deleted{\unsure{u}}\epsilon(x)}_y-
\epsilon\pi\underbrace{(\epsilon(x))}=\lambda y-\underbrace{\epsilon\pi(y)}_{\lambda y}=0
\]
d'autre par sa restriction à $E^0$ est $\lambda\,\mathrm{id}_{\unsure{E^0|E}}$
\[
u\mid E^0=\lambda\,\mathrm{id}_{E^0}\qquad
\begin{aligned}&\text{i.e. }u(x)=\lambda x\\&\text{si }\unsure{\alpha|x}\in E^0\end{aligned}
\]

% L12–19
C'est la version duale du $\lambda$-scindage.
On récupère $\pi$ par $u$, par la
\unsure{forme|form.|formule}
\[
\pi(\epsilon(x))=\lambda x-u(x)
\]
en notant que le deuxième \unsure{mbre|membre} se
factorise bien par $F$, car il est nul
sur $E^0$. La donnée de $\pi$ ou de $u$
est équivalente. Mais c'est $u$ qui
sera utilisé pour la théorie des $*_{i}$.
% archive sheet number: 174
\scanpage{176}
% L1–6
On se fixe un $\lambda\in k$, et on
travaille avec les extensions $\lambda$-scindées.
\deleted{On appelle} \unsure{se donné|se donne|Si on se donne} \unsure{de|des} telles
\diagram{\begin{tikzcd}[ampersand replacement=\&,column sep=small]
0\arrow[r,dash]\&E_i^0\arrow[r,dash,"\mathrm{inj}_i"]\&
 E_i\arrow[r,dash]\arrow[l,bend left=30,"u_i"']\& F_i\arrow[r,dash]\&0
\end{tikzcd}}
\[
u_i\circ\mathrm{inj}_i=\lambda\unsure{\xi|\epsilon|{}}\,\mathrm{id}_{E_i^0}
\qquad \lambda\ \text{\uuline{indép. de} }i
\]
et un module $B$\unsure{,} sur $k$. \deleted{Une}

% L7–12
On appelle $\lambda$-accouplement à valeurs
dans $B$, une couple
\[
\begin{aligned}
\prod E_i &\mathrel{\mathop{\text{---}}\limits^{f}} B\\
\prod E_j^0 &\mathrel{\mathop{\text{---}}\limits^{f^0}} B
\end{aligned}
\]
d'applications \deleted{\unsure{bi}} \inserted{multi}linéaires, satisfaisant
\[
f(x_1,\ldots,x_i+\alpha_i,\ldots x_n)
=f(x_1,\ldots x_i,\ldots x_n)+f^0(u_1x_1),\ldots,\alpha_i,\ldots,u_n(x_n))
\]
% L13–19
Prenant les $x_i$ dans $E_i^0$, on trouve
\[
f(\alpha_1,\ldots,\alpha_i,\ldots\alpha_n)=\lambda^{n-1}f^0(\alpha_1,\ldots,\alpha_i,\ldots\alpha_n)
\]
(on a pris $x_j=\alpha_j$ pour $j=1,\ldots\hat i,\ldots n$, et
$x_i=0$). On a une solution du
pb universel, noté $\mathop{*}\limits_{i,\lambda}E_i$ ($E_1*_\lambda E_2$ etc)
l'indice $\lambda$ rappelant que c'est relatif à un
choix de $\lambda$-trivialisations des extensions
envisagées. L'\unsure{espace|expression} $\mathop{*}\limits_{\unsure{i|j},\lambda}E_i$ se construit
% archive sheet number: 175
\scanpage{177}
% L1–4
comme quotient de
\[
\bigotimes_i E_i \oplus \bigotimes_i E_i^0
\]
et on a une suite exacte
\[
\bigotimes E_i^0 \longrightarrow \mathop{\star}\limits_{i,\lambda} E_i \unsure{\text{---}|\longrightarrow} \bigotimes \deleted{\unsure{E}}F_i\deleted{\unsure{S|s}}\longrightarrow 0
\]
% L5–8
qui correspond au fait que les systèmes $(f,f^0)$ pour lesquels $f^0=0$, \unsure{correspondent|s'identifient} aux $f$ qui \deleted{\unsure{q|se}} factorisent par les $F_i$.

% L9–13
L'application $\bigotimes E_i^0 \mathrel{\overset{\mathrm{inj}}{\unsure{\text{---}|\longrightarrow}}} \mathop{\star}\limits_{i,\lambda} E_i$ n'est sans doute pas tjrs injective, mais je dis que le choix des scindages des extensions $E_i$ définit un inverse à gauche de $\mathrm{inj}$.

% L14–19
En effet, écrivons
\[
E_i \simeq E_i^0\oplus F_i
\]
la donnée de $u_i$, déjà connue sur $E_i^0$ comme $\lambda E_i^0$, équivaut alors à la donnée de $\deleted{\unsure{t|b}}b_i:F_i \to E_i^0$
% archive sheet number: 176
